\section{优先池总览}

16只优先股全部具备Q1结构化财务并建立公司级熊／基准／牛模型；15只归档最新原始券商PDF，浙江东方研报接口失败但仍以官方预告、Q1财务和PB-ROE/盈利桥建立House区间。对浙江东方已检查结构化研报接口、官方H1预告、Q1财务和公开检索；无法取得可核验的当前原始券商PDF，因此Street权重为0，但House模型仍可由公告分母复算。外部目标与AStock House目标严格分列。

\section{静默行业验证：1只}

\begin{exhibitbox}[Exhibit 7：静默行业验证]
\scriptsize
\begin{tabularx}{\textwidth}{L{1.0cm}L{1.45cm}R{0.85cm}R{0.85cm}R{0.95cm}R{0.85cm}R{0.85cm}R{0.85cm}X}
\toprule
\textbf{代码} & \textbf{公司} & \textbf{现价} & \textbf{下沿} & \textbf{概率目标} & \textbf{上沿} & \textbf{空间} & \textbf{H1扣非} & \textbf{动作} \\
\midrule
601225 & 陕西煤业 & 22.92 & 13.60 & 19.76 & 26.40 & -14\% & 97.5 & 估值已满／仅观察 \\
\bottomrule
\end{tabularx}
\normalsize
\end{exhibitbox}

\paragraph{陕西煤业（601225）}
现价22.92元，一年位置55.3\%；H1归母／扣非净利中值114.6/97.5亿元；Q1归母净利／经营现金流42.1/96.8亿元。最新研报为国信证券 2026-04-28《煤价下行拖累业绩，2026年业绩弹性可期》，2026E EPS 2.08元、PE 12.2倍，状态为current。估值采用周期调整PE与股息底部交叉验证，熊／基准／牛为13.60/20.80/26.40元，概率目标19.76元，相对现价-13.8\%。
\textbf{催化}：煤价企稳、2026E EPS高于2.08元且分红稳定；\textbf{失效}：煤价下跌、扣非利润走弱或分红下降。证据等级medium-high，外部目标权重0\%；未引用外部目标时，概率目标明确为AStock House模型。


\section{低位盈利优先池：11只}

\begin{exhibitbox}[Exhibit 8：低位盈利优先池]
\scriptsize
\begin{tabularx}{\textwidth}{L{1.0cm}L{1.45cm}R{0.85cm}R{0.85cm}R{0.95cm}R{0.85cm}R{0.85cm}R{0.85cm}X}
\toprule
\textbf{代码} & \textbf{公司} & \textbf{现价} & \textbf{下沿} & \textbf{概率目标} & \textbf{上沿} & \textbf{空间} & \textbf{H1扣非} & \textbf{动作} \\
\midrule
002379 & 宏桥控股 & 14.45 & 12.50 & 22.42 & 32.50 & +55\% & 154.0 & 高空间模型候选／先验证 \\
600346 & 恒力石化 & 15.79 & 11.20 & 17.18 & 24.60 & +9\% & 53.3 & 市场支持观察／等安全边际 \\
002738 & 中矿资源 & 51.80 & 35.00 & 57.53 & 88.00 & +11\% & 11.0 & 市场支持观察／等安全边际 \\
002048 & 宁波华翔 & 23.04 & 14.40 & 23.92 & 38.00 & +4\% & 5.3 & 市场支持观察／等安全边际 \\
002532 & 天山铝业 & 11.17 & 10.50 & 16.54 & 24.20 & +48\% & 40.9 & 高空间模型候选／先验证 \\
600595 & 中孚实业 & 5.94 & 5.25 & 8.61 & 12.65 & +45\% & 19.2 & 高空间模型候选／先验证 \\
601360 & 三六零 & 9.06 & 2.14 & 3.60 & 5.71 & -60\% & 1.6 & 高估值风险／勿追 \\
600918 & 中泰证券 & 5.30 & 5.10 & 7.42 & 10.48 & +40\% & 17.0 & 高空间模型候选／先验证 \\
300014 & 亿纬锂能 & 54.94 & 41.60 & 66.13 & 98.80 & +20\% & 25.2 & 选择性回撤／业绩验证 \\
000987 & 越秀资本 & 8.09 & 5.52 & 7.76 & 10.64 & -4\% & 23.3 & 估值已满／仅观察 \\
600120 & 浙江东方 & 4.84 & 3.23 & 4.75 & 6.72 & -2\% & 8.4 & 估值已满／仅观察 \\
\bottomrule
\end{tabularx}
\normalsize
\end{exhibitbox}

\paragraph{宏桥控股（002379）}
现价14.45元，一年位置9.3\%；H1归母／扣非净利中值155.0/154.0亿元；Q1归母净利／经营现金流67.6/92.4亿元。最新研报为招银国际 2026-05-13《高盈利弹性，高分红比例；首次覆盖给予买入评级》，2026E EPS ---元、PE 10.5倍，状态为current。估值采用扣非EPS周期PE，并以原始券商目标交叉验证，熊／基准／牛为12.50/23.65/32.50元，概率目标22.42元，相对现价+55.2\%。
\textbf{催化}：铝价成本差与高分红延续，H2扣非利润保持强势；\textbf{失效}：铝价成本差收窄，或分红及现金转化恶化。证据等级high，外部目标权重5\%；未引用外部目标时，概率目标明确为AStock House模型。

\paragraph{恒力石化（600346）}
现价15.79元，一年位置17.7\%；H1归母／扣非净利中值72.0/53.3亿元；Q1归母净利／经营现金流39.1/45.5亿元。最新研报为中邮证券 2026-07-10《2026H1业绩高增，“大化工”平台优势显著》，2026E EPS 1.78元、PE 13.3倍，状态为current。估值采用预告后正常化PE，并对一次性收益折价，熊／基准／牛为11.20/17.80/24.60元，概率目标17.18元，相对现价+8.8\%。
\textbf{催化}：炼化价差稳定且H2经营现金流为正；\textbf{失效}：一次性收益占比上升、价差收窄或H2利润低于基准桥。证据等级medium-high，外部目标权重0\%；未引用外部目标时，概率目标明确为AStock House模型。

\paragraph{中矿资源（002738）}
现价51.80元，一年位置32.3\%；H1归母／扣非净利中值11.5/11.0亿元；Q1归母净利／经营现金流5.1/0.2亿元。最新研报为联储证券 2026-05-07《中矿资源一季报点评：锂盐价格同比高增，推动业绩强势释放》，2026E EPS 3.27元、PE 25.4倍，状态为current。估值采用锂周期PE，并以资源产量和现金流验证，熊／基准／牛为35.00/58.86/88.00元，概率目标57.53元，相对现价+11.1\%。
\textbf{催化}：锂价、资源产量与H2现金转化共同支撑EPS；\textbf{失效}：锂价反转、产量不及预期或经营现金流持续偏弱。证据等级medium-high，外部目标权重0\%；未引用外部目标时，概率目标明确为AStock House模型。

\paragraph{宁波华翔（002048）}
现价23.04元，一年位置25.6\%；H1归母／扣非净利中值6.5/5.3亿元；Q1归母净利／经营现金流2.7/1.9亿元。最新研报为国金证券 2025-10-30《Q3业绩超预期，机器人布局加速》，2026E EPS 1.98元、PE 17.3倍，状态为stale。估值采用汽车零部件正常化PE，并对陈旧预测折价，熊／基准／牛为14.40/24.00/38.00元，概率目标23.92元，相对现价+3.8\%。
\textbf{催化}：欧洲业务修复和机器人增量推升扣非利润；\textbf{失效}：H2扣非利润低于H1节奏或海外亏损再次扩大。证据等级medium，外部目标权重0\%；未引用外部目标时，概率目标明确为AStock House模型。

\paragraph{天山铝业（002532）}
现价11.17元，一年位置26.5\%；H1归母／扣非净利中值42.0/40.9亿元；Q1归母净利／经营现金流22.2/23.4亿元。最新研报为国金证券 2026-04-29《稀缺产能释放，全年开局表现亮眼》，2026E EPS 1.92元、PE 8.8倍，状态为current。估值采用电解铝周期PE，并验证产能释放与分红，熊／基准／牛为10.50/17.10/24.20元，概率目标16.54元，相对现价+48.1\%。
\textbf{催化}：低成本产能释放且分红支撑约1.9元EPS桥；\textbf{失效}：铝价差、产量或现金流显著走弱。证据等级medium-high，外部目标权重0\%；未引用外部目标时，概率目标明确为AStock House模型。

\paragraph{中孚实业（600595）}
现价5.94元，一年位置28.4\%；H1归母／扣非净利中值18.8/19.2亿元；Q1归母净利／经营现金流8.2/-0.8亿元。最新研报为国信证券 2026-04-28《电解铝和加工业务单位盈利高增，2025年分红比例43\%》，2026E EPS 1.00元、PE 8.5倍，状态为current。估值采用周期PE并计入现金流折价，熊／基准／牛为5.25/9.00/12.65元，概率目标8.61元，相对现价+45.0\%。
\textbf{催化}：电解铝与加工业务单位利润稳定、经营现金流转正；\textbf{失效}：经营现金流继续为负或单位盈利快速正常化。证据等级medium-high，外部目标权重0\%；未引用外部目标时，概率目标明确为AStock House模型。

\paragraph{三六零（601360）}
现价9.06元，一年位置12.9\%；H1归母／扣非净利中值2.2/1.6亿元；Q1归母净利／经营现金流1.1/-0.6亿元。最新研报为信达证券 2023-04-25《安全业务打造二次成长曲线，积极布局AI展望行业蓝海》，2026E EPS ---元、PE ---倍，状态为stale。估值采用收入倍数与盈利路径交叉验证，熊／基准／牛为2.14/3.64/5.71元，概率目标3.60元，相对现价-60.3\%。
\textbf{催化}：安全与AI收入提速、经营现金流转正；\textbf{失效}：收入停滞、现金流持续为负或AI商业化贡献不显著。证据等级medium，外部目标权重0\%；未引用外部目标时，概率目标明确为AStock House模型。

\paragraph{中泰证券（600918）}
现价5.30元，一年位置8.8\%；H1归母／扣非净利中值17.5/17.0亿元；Q1归母净利／经营现金流4.7/49.2亿元。最新研报为太平洋 2025-10-31《中泰证券2025年中报点评：财富管理表现出色，自营投资提供弹性》，2026E EPS 0.24元、PE 28.8倍，状态为stale。估值采用PB-ROE与盈利情景混合估值，熊／基准／牛为5.10/7.59/10.48元，概率目标7.42元，相对现价+40.0\%。
\textbf{催化}：资本市场活跃度和财富管理收入支撑H2利润；\textbf{失效}：自营收益反转，或ROE、手续费收入低于基准。证据等级medium，外部目标权重0\%；未引用外部目标时，概率目标明确为AStock House模型。

\paragraph{亿纬锂能（300014）}
现价54.94元，一年位置24.4\%；H1归母／扣非净利中值32.5/25.2亿元；Q1归母净利／经营现金流14.5/-3.7亿元。最新研报为群益证券 2026-06-26《公司26H1净利润翻倍增长，储能业务加快扩张，建议“买进”》，2026E EPS 3.11元、PE 20.0倍，状态为current。估值采用预告后PE，并验证储能出货与现金流，熊／基准／牛为41.60/65.35/98.80元，概率目标66.13元，相对现价+20.4\%。
\textbf{催化}：储能出货、毛利率和H2现金转化验证EPS；\textbf{失效}：经营现金流持续为负，或储能价格与毛利低于预期。证据等级medium-high，外部目标权重10\%；未引用外部目标时，概率目标明确为AStock House模型。

\paragraph{越秀资本（000987）}
现价8.09元，一年位置33.0\%；H1归母／扣非净利中值28.8/23.3亿元；Q1归母净利／经营现金流13.0/39.1亿元。最新研报为太平洋 2019-08-19《越秀金控中报点评：证券业务业绩波动剧烈，业务转型进行中》，2026E EPS ---元、PE ---倍，状态为stale。估值采用PB-ROE与盈利分部混合估值，熊／基准／牛为5.52/7.96/10.64元，概率目标7.76元，相对现价-4.1\%。
\textbf{催化}：投资与租赁盈利强劲、账面价值稳定增长；\textbf{失效}：投资收益反转或ROE低于基准。证据等级medium，外部目标权重0\%；未引用外部目标时，概率目标明确为AStock House模型。

\paragraph{浙江东方（600120）}
现价4.84元，一年位置17.6\%；H1归母／扣非净利中值8.7/8.4亿元；Q1归母净利／经营现金流1.5/0.6亿元。已检查结构化研报接口、官方公告与公开检索；当前原始券商PDF未获得，外部目标权重为0，House分母来自H1预告与Q1财务。估值采用PB-ROE与盈利分部混合估值，熊／基准／牛为3.23/4.88/6.72元，概率目标4.75元，相对现价-1.9\%。
\textbf{催化}：金融持股与经营利润转化为持续ROE；\textbf{失效}：H2利润减速或金融资产收益反转。证据等级medium-low，外部目标权重0\%；未引用外部目标时，概率目标明确为AStock House模型。


\section{已启动仍有空间候选：4只}

\begin{exhibitbox}[Exhibit 9：已启动仍有空间候选]
\scriptsize
\begin{tabularx}{\textwidth}{L{1.0cm}L{1.45cm}R{0.85cm}R{0.85cm}R{0.95cm}R{0.85cm}R{0.85cm}R{0.85cm}X}
\toprule
\textbf{代码} & \textbf{公司} & \textbf{现价} & \textbf{下沿} & \textbf{概率目标} & \textbf{上沿} & \textbf{空间} & \textbf{H1扣非} & \textbf{动作} \\
\midrule
000703 & 恒逸石化 & 13.73 & 11.40 & 16.26 & 22.50 & +18\% & 57.2 & 选择性回撤／业绩验证 \\
000301 & 东方盛虹 & 12.55 & 10.50 & 18.54 & 29.70 & +48\% & 44.2 & 高空间模型候选／先验证 \\
002414 & 高德红外 & 14.64 & 11.25 & 20.59 & 33.75 & +41\% & 13.2 & 高空间模型候选／先验证 \\
002558 & 巨人网络 & 29.56 & 19.00 & 31.45 & 48.60 & +6\% & 23.0 & 市场支持观察／等安全边际 \\
\bottomrule
\end{tabularx}
\normalsize
\end{exhibitbox}

\paragraph{恒逸石化（000703）}
现价13.73元，一年位置64.1\%；H1归母／扣非净利中值57.5/57.2亿元；Q1归母净利／经营现金流19.9/-25.7亿元。最新研报为国信证券 2026-07-05《文莱炼化项目盈利提升，公司上半年业绩大幅增长》，2026E EPS 2.20元、PE 6.9倍，状态为current。估值采用预告后炼化周期PE，并使用券商合理价值区间，熊／基准／牛为11.40/16.50/22.50元，概率目标16.26元，相对现价+18.4\%。
\textbf{催化}：文莱炼化盈利和H2现金流验证预告后区间；\textbf{失效}：炼化价差收窄或经营现金流持续为负。证据等级high，外部目标权重10\%；未引用外部目标时，概率目标明确为AStock House模型。

\paragraph{东方盛虹（000301）}
现价12.55元，一年位置67.4\%；H1归母／扣非净利中值46.0/44.2亿元；Q1归母净利／经营现金流14.3/35.3亿元。最新研报为中邮证券 2026-05-06《石化新材料销量稳步增长，高端新材料矩阵持续丰富》，2026E EPS 0.47元、PE 28.3倍，状态为current。估值采用石化与新材料结构的预告后重置PE，熊／基准／牛为10.50/18.90/29.70元，概率目标18.54元，相对现价+47.7\%。
\textbf{催化}：H2新材料占比和正现金流支撑新盈利分母；\textbf{失效}：价差正常化或H2扣非利润低于基准桥。证据等级medium-high，外部目标权重0\%；未引用外部目标时，概率目标明确为AStock House模型。

\paragraph{高德红外（002414）}
现价14.64元，一年位置54.8\%；H1归母／扣非净利中值13.6/13.2亿元；Q1归母净利／经营现金流3.7/-1.1亿元。最新研报为太平洋 2025-11-06《合同负债大幅增长，看好全年业绩表现》，2026E EPS 0.18元、PE 73.5倍，状态为stale。估值采用军工预告后重置PE，并验证订单与现金流，熊／基准／牛为11.25/21.00/33.75元，概率目标20.59元，相对现价+40.6\%。
\textbf{催化}：订单、合同负债和H2现金流支撑预告后重置；\textbf{失效}：订单转化减速或经营现金流持续为负。证据等级medium，外部目标权重5\%；未引用外部目标时，概率目标明确为AStock House模型。

\paragraph{巨人网络（002558）}
现价29.56元，一年位置24.2\%；H1归母／扣非净利中值21.0/23.0亿元；Q1归母净利／经营现金流10.8/13.2亿元。最新研报为东吴证券 2026-07-09《26H1业绩预告点评：业绩符合预期，出海增量可期》，2026E EPS 2.29元、PE 12.2倍，状态为current。估值采用预告后游戏PE，并验证新品与出海，熊／基准／牛为19.00/32.06/48.60元，概率目标31.45元，相对现价+6.4\%。
\textbf{催化}：新品和海外收入支撑H2扣非利润；\textbf{失效}：产品管线或商业化低于预期，H2利润低于基准桥。证据等级medium-high，外部目标权重0\%；未引用外部目标时，概率目标明确为AStock House模型。

\begin{keyinsight}[优先池的真正排序]
按概率目标空间排序，宏桥控股、天山铝业、东方盛虹、中孚实业、高德红外和中泰证券居前；但它们分别受铝价差、产能释放、石化价差、负现金流、订单兑现和资本市场波动约束。三六零、陕西煤业、越秀资本与浙江东方的概率目标不高于现价，说明低位置或高同比并不自动等于低估。
\end{keyinsight}
